\chapter{第 2 段階の詳細計画：アセンブル・校正・実現性の検証}\label{ch:phase2}

第~\ref{ch:conclusion}章で示した第 2 段階の構想を、第 1 段階の結果を使って具体化する。
目標は、\textbf{組み合わせた形状（アセンブル）について最適な動作の校正値を定義し、要求した動きや性能を実現できるかを検証する}ことである。

\section{アセンブルの部品表}

第 1 段階で評価した形状の要素を、組み合わせ可能な「部品」として整理する。

\begin{longtable}{p{0.18\linewidth}p{0.36\linewidth}p{0.38\linewidth}}
\toprule
部品の種類 & 選択肢 & 主なパラメータと評価した章 \\ \midrule \endhead
下半身 & 二本脚／戦車型（履帯）／車輪 & 支持範囲、重心の高さ、脚の質量（第~\ref{ch:tank}, \ref{ch:legs}, \ref{ch:steps}章） \\
腰・上体 & 固定／前後傾 1 軸／2 軸 & 傾きの範囲・速度・加速度、上体の質量（第~\ref{ch:legs}, \ref{ch:steps}章） \\
肩 & 2 軸ジンバル／3 軸／3 軸＋上肢帯 & 可動域（水平内転、内旋・外旋、前方突出）、特異点の位置（第~\ref{ch:3d}, \ref{ch:shoulder}章） \\
腕 & 1 リンク／2 リンク（肘あり） & 長さ、質量、関節の速度・加速度の上限（第~\ref{ch:tank}, \ref{ch:3d}章） \\
胸・装甲 & 薄い／標準／厚い & 胸の奥行（第~\ref{ch:shoulder}章） \\
質量配分 & 上半身重視／脚重視 & 手足の質量比（第~\ref{ch:planar}, \ref{ch:legs}章） \\
姿勢制御の補助 & なし／RW／CMG／スラスタ & 容量、トルク、推進剤（第~\ref{ch:rw}, \ref{ch:ai}章） \\
アクチュエータ & 電動／油圧 & トルク密度、帯域、遅れ（今後の課題） \\
大きさ・環境 & 人間大〜18 m 級、地球・火星・月・宇宙 & 長さの倍率 $s$、重力 $g$（第~\ref{ch:dyn}, \ref{ch:tank}, \ref{ch:steps}章） \\
\bottomrule
\end{longtable}

\section{評価指標（性能表）}

組み合わせた機体ごとに、第 1 段階のシミュレーションを自動で実行し、次の性能表を作る。

\begin{longtable}{p{0.30\linewidth}p{0.20\linewidth}p{0.40\linewidth}}
\toprule
指標 & 単位 & 求め方 \\ \midrule \endhead
1 周あたりの回転（宇宙） & ° & 曲率の面積分の上界（第~\ref{ch:planar}章）、3 次元は円錐運動（第~\ref{ch:3d}章） \\
90° 旋回時間（宇宙） & s & トルク上限から $f_{\max}$（第~\ref{ch:dyn}章） \\
角運動量の容量 & kg\,m$^2$/s & $h_{\mathrm{AMBAC}}$, $h_{\mathrm{RW}}$（第~\ref{ch:rw}章） \\
推進剤（宇宙、時間指定の旋回） & kg & 線形計画（第~\ref{ch:ai}章） \\
荷物運びの最短時間（地上） & s & 転倒・滑り制約下の軌道最適化（第~\ref{ch:tank}, \ref{ch:legs}章） \\
定格荷重マップ（地上） & kg & 方向ごとの静的な最大荷重（第~\ref{ch:3d}章） \\
旋回の安全な速さ（地上） & s & 3 次元の旋回最適化（第~\ref{ch:3d}章） \\
意図しない旋回（低重力） & °/往復 & スティックスリップ（第~\ref{ch:tank}章） \\
越えられる段差 & m & 静力学（第~\ref{ch:steps}章） \\
人間らしい姿勢の可否 & 可／不可、必要可動域 & 逆運動学（第~\ref{ch:shoulder}章） \\
\bottomrule
\end{longtable}

\section{校正値の定義}

シミュレーションと実機の差を埋めるために、次の値を測定・推定して与える。

\begin{longtable}{p{0.22\linewidth}p{0.36\linewidth}p{0.34\linewidth}}
\toprule
校正値 & 測定・推定の方法 & 影響の大きい指標 \\ \midrule \endhead
慣性パラメータ（各リンクの質量・重心・慣性） & CAD の値＋AMBAC ループによる自己同定（次節） & 1 周あたりの回転、旋回時間、ZMP \\
摩擦係数 & 斜面の滑り試験、引っ張り試験 & 滑り、段差、スティックスリップ \\
接地範囲 & 履帯・足裏の寸法と、変形を考慮した有効範囲 & 転倒、定格荷重 \\
関節の可動域 & 実機で各関節を端まで動かして測る（ひねりは測り方を明記） & 姿勢の可否、ループの上界 \\
アクチュエータ特性 & トルク・速度の上限、帯域、遅れ（ステップ応答） & 旋回時間、最短動作時間 \\
人の動きのデータ & モーションキャプチャー（目標の形として） & 人間らしさの評価、学習の構造 \\
\bottomrule
\end{longtable}

\section{AMBAC ループによる慣性パラメータの自己同定}

宇宙や空中（外力がない状態）では、機体が自分の手足をループ状に動かし、胴体の回転を測れば、慣性パラメータを推定できる可能性がある。
外部の計測器を使わずに済むので、実機の「自己校正」として魅力的である。
その実現性を、人間大の平面モデルで調べた（図~\ref{fig:D14_self_identification.png}）。

\paragraph{設定}
推定するパラメータは、腕の質量、脚の質量、胴体の慣性（倍率）の 3 つとした（総質量は既知とする）。
曲率マップの異なる場所に置いた 6 種類の円ループ（L1〜L6、L1 は Sim03 の最適ループ）を、それぞれ 10 周ずつ回す。
姿勢の計測誤差を 1 回あたり 0.1° とし、パラメータの推定誤差を、感度（ヤコビ行列 $J$）から最小二乗の共分散 $\sigma^2(J^\top J)^{-1}/N$ で見積もった。

\paragraph{結果}
\begin{longtable}{p{0.40\linewidth}rrr}
\toprule
使う情報 & 腕の質量 & 脚の質量 & 胴体の慣性 \\ \midrule \endhead
1 周ごとの正味の回転だけ（6 個の値） & 29.9\% & 223.8\% & 4.1\% \\
周回中の姿勢の波形（各ループ 36 点） & 3.2\% & 10.1\% & 0.9\% \\
\bottomrule
\end{longtable}

\begin{itemize}
\item \textbf{正味の回転だけでは、脚の質量はほとんど推定できない}（推定誤差 224\%）。総質量を固定すると、脚を重くした分だけ胴体が軽くなり、
  正味の回転への影響が打ち消し合うためである。胴体の慣性は、正味の回転への影響が大きいのでよく推定できる。
\item \textbf{周回中の姿勢の波形まで使うと、推定誤差は 5〜20 分の 1 になり、3 つとも推定できる}。
  波形には、各形状での $\Gamma$（接続）の情報が含まれるためである。
\end{itemize}

\DFIG{D14_self_identification.png}{AMBAC ループによる自己同定。左：各ループの回転の、パラメータへの感度（1\% あたり）。右：姿勢の計測誤差と推定誤差（実線：正味の回転だけ、破線：周回中の波形も使う）。}

\begin{keypoint}
\textbf{校正の指針}：AMBAC ループによる自己同定では、ループ 1 周の正味の回転だけでなく、周回中の姿勢の時間変化（IMU の波形）を記録して使う。
曲率マップの異なる場所に複数のループを置くと、パラメータどうしの区別がつきやすくなる。
\end{keypoint}

\section{実現性の検証の手順}

\begin{algorithm}[htbp]
\caption{アセンブル型の設計・校正・検証}
\begin{algorithmic}[1]
\Require 要求（動き・性能の目標値）、部品の候補
\State 要求から、必要な評価指標と合格の基準を決める
\For{各部品の組み合わせ（アセンブル）}
  \State シミュレーションで性能表を作る（第 1 段階の評価を自動実行）
  \State 合格基準を満たさなければ、次の組み合わせへ
\EndFor
\State 候補の機体について校正値を測る（自己同定、摩擦試験、可動域の測定など）
\State 校正値を入れて性能表を作り直し、要求を満たすか検証する
\If{満たさない}
  \State 制御（ゲイン、軌道最適化、学習方策の制約）を調整する。それでも満たさなければ部品を選び直す
\EndIf
\end{algorithmic}
\end{algorithm}

\section{例：要求から部品を選ぶ}

第 1 段階の結果だけを使って、次の要求を持つ人間大の機体の部品を選んでみる。

\begin{quote}
要求：(1) 腕組みができる。(2) 20 cm の段差を越える。(3) 月面で 10 kg の荷物を運ぶ作業をする。(4) 作業中に機体が意図せず回らない。
\end{quote}

\begin{longtable}{p{0.20\linewidth}p{0.36\linewidth}p{0.36\linewidth}}
\toprule
要求 & 第 1 段階の結果 & 選ぶ部品・運用 \\ \midrule \endhead
(1) 腕組み & 2 軸の肩では不可能。3 軸の肩なら薄い胸で人の限界ぎりぎり。上肢帯があれば胸の厚さによらず可能（第~\ref{ch:shoulder}章） & 3 軸の肩＋上肢帯（前方突出 20° 程度） \\
(2) 20 cm の段差 & 戦車型は 25 cm まで可能（形の限界）。二本脚は上体の前傾ありで 60 cm まで、前傾なしだと 18 cm で不可（第~\ref{ch:steps}章） & 戦車型でも二本脚（腰の関節つき）でもよい \\
(3) 月面の荷物運び & 地球より 2.5 倍遅く動かす必要がある（フルード則）。軌道を最適化すると約 2 s（第~\ref{ch:tank}章） & 軌道最適化を使う。二本脚なら腰の戦略が必須（第~\ref{ch:legs}章） \\
(4) 意図しない旋回 & 月面では 2.1 s より速い腕の振りで機体が滑って回る（第~\ref{ch:tank}章） & 腕の水平の振りを 2 s 以上かける、または摩擦の大きい接地面を使う \\
\bottomrule
\end{longtable}

この例では、「3 軸の肩＋上肢帯」と「腰の関節」が必須で、下半身は段差の要求だけなら戦車型でも二本脚でもよい。
転倒のしにくさとコスト（第~\ref{ch:legs}章）を重視すれば戦車型＋腰、より高い段差や不整地を重視すれば二本脚になる。
第 2 段階では、このような選択を、校正値を入れたシミュレーションで定量的に検証する。

\section{第 2 段階で新たに必要になるもの}

\begin{itemize}
\item \textbf{評価の自動化}：部品の組み合わせを入力として、第 1 段階の 10 本のシミュレーションをまとめて実行し、性能表を出力する仕組み。
\item \textbf{アクチュエータの簡易モデル}：帯域・遅れ・トルク速度制限・コンプライアンスを持つモデル。油圧と電動の比較。
\item \textbf{自己同定のアルゴリズム}：周回中の姿勢の波形から慣性パラメータを推定する最小二乗法（またはカルマンフィルタ）。
\item \textbf{人の動きのデータとの接続}：モーションキャプチャーのデータを、形状の異なる機体に移す方法（ひねりの扱いを含む）。
\item \textbf{動的な動作}：踏み出し、勢いを使った段差の乗り越え。
\end{itemize}
